\chapter{Malware}

\section{Introduzione}

\begin{defbox}[Malware]
Il malware (malicious software) è un programma inserito in un sistema con l'intento di comprometterne la confidenzialità, l'integrità o la disponibilità dei dati.
\end{defbox}

\subsection{Classificazione e Propagazione}

Il malware può essere classificato in due grandi categorie: in base ai \textbf{meccanismi di propagazione} e alle \textbf{azioni eseguite} (o payload) una volta raggiunto il bersaglio.

\begin{notebox}[Meccanismi di propagazione]
La prima categoria di propagazione del malware riguarda frammenti di \textit{software parassita} che si attaccano a contenuti eseguibili esistenti. Il frammento può essere:
\begin{itemize}
    \item Codice macchina che infetta un'applicazione, un programma di utilità o il codice utilizzato per l'avvio del sistema.
    \item Codice di scripting, utilizzato in file di dati come documenti Word, fogli di calcolo Excel o documenti PDF.
\end{itemize}
\end{notebox}

\section{Tipi di Malware}

Il malware comprende una vasta e variegata serie di software dannosi progettati per compromettere la sicurezza, la privacy o l'integrità dei sistemi informatici. Di seguito sono presentati i principali tipi di malware, classificati per caratteristiche e modalità di azione.

\begin{warnbox}[Advanced Persistent Threat (APT)]
Un APT è una minaccia avanzata e persistente che utilizza diverse forme di malware per attaccare un obiettivo specifico in modo continuo e discreto per un lungo periodo. Spesso orchestrati da gruppi di hacker organizzati (o nazioni), mirano a ottenere l'accesso prolungato a una rete per rubare dati sensibili o sabotarne le operazioni.
\end{warnbox}

\begin{defbox}[Tipologie Comuni di Malware]
\begin{description}
    \item[Adware] Software che visualizza o scarica pubblicità indesiderata (pop-up, reindirizzamenti). Oltre a compromettere l'esperienza utente, può raccogliere informazioni di navigazione ed essere difficile da rimuovere.
    \item[Exploit Kit] Strumenti che automatizzano la creazione e la distribuzione di nuovo malware sfruttando vulnerabilità software note. Offrono meccanismi di propagazione sofisticati. Esempi noti: \textit{Zeus}, \textit{Blackhole}, \textit{Sakura} e \textit{Phoenix}.
    \item[Backdoor] Meccanismi che consentono di aggirare i normali controlli di sicurezza, garantendo un accesso non autorizzato a un sistema o a un'applicazione.
    \item[Exploits] Frammenti di codice progettati per sfruttare specifiche vulnerabilità nei software per eseguire codice non autorizzato.
    \item[Flooders (DoS Client)] Strumenti utilizzati per generare un grande volume di traffico al fine di inondare e sovraccaricare un sistema (attacchi DoS).
    \item[Keyloggers] Malware che monitorano e registrano ogni tasto premuto, fornendo agli attaccanti password o numeri di carta di credito.
    \item[Logic Bomb] Malware programmati per rimanere inattivi fino al verificarsi di una condizione predefinita (es. una data). Una volta attivate, eseguono azioni dannose.
    \item[Macro Virus] Virus che utilizzano codice o script incorporati in documenti (es. Word o Excel), attivandosi all'apertura del file.
    \item[Mobile Code] Script o software trasferibili ed eseguibili su diverse piattaforme senza modifiche (es. JavaScript, ActiveX).
    \item[Rootkit] Strumenti progettati per ottenere e mantenere il controllo root (amministratore) su un sistema senza essere rilevati.
    \item[Spyware] Raccoglie informazioni dal computer dell'utente a sua insaputa (digitazione, dati a schermo, file di sistema) e le trasmette a terzi.
    \item[Trojan] Programmi che appaiono utili ma nascondono funzioni dannose. Eludono la sicurezza sfruttando autorizzazioni legittime per poi installare altro malware.
    \item[Virus] Malware che tenta di replicarsi inserendosi in altri file o programmi eseguibili.
    \item[Worm] Programma che si replica autonomamente e si diffonde in rete sfruttando vulnerabilità software.
    \item[Zombie / Bot] Programma che, una volta installato, può essere utilizzato remotamente per lanciare attacchi coordinati (spesso in una \textit{botnet}).
\end{description}
\end{defbox}

\section{Virus}
Un virus informatico è un software che infetta altri programmi modificandoli con l'iniezione di codice che gli consente di replicarsi.
Ogni volta che il computer infetto entra in contatto con codice non infetto, il virus si trasferisce diffondendosi da macchina a macchina. La diffusione può avvenire tramite supporti fisici (chiavette USB) o tramite reti, sfruttando documenti o servizi condivisi.

\subsection{Struttura Generale e Fasi}

\begin{notebox}[Struttura di un Virus]
Un virus è costituito generalmente da tre parti:
\begin{enumerate}
    \item \textbf{Meccanismo di infezione} (Vettore di infezione): il mezzo con cui il virus si propaga.
    \item \textbf{Trigger} (Bomba logica): l'evento che determina l'attivazione o la consegna del payload.
    \item \textbf{Payload}: ciò che il virus fa effettivamente oltre a diffondersi (es. danni al sistema o attività benigne ma fastidiose).
\end{enumerate}
Una volta avviato, il virus cerca file eseguibili non infetti per replicarsi, dopodiché esegue il suo payload se le condizioni del trigger sono soddisfatte, per poi restituire il controllo al programma originale.
\end{notebox}

Spesso un virus comprime il file eseguibile infetto per mascherare l'aumento di dimensione dovuto all'iniezione di codice, rimanendo nascosto durante l'esecuzione del programma.

\begin{center}
\begin{tikzpicture}[node distance=1.5cm, auto, thick, scale=0.8, transform shape]
\node (dorm) [draw, rounded corners, fill=blue!10, text width=2.5cm, align=center] {\textbf{Fase Dormiente}\\Inattivo fino a trigger};
\node (prop) [draw, rounded corners, fill=green!10, right=of dorm, text width=2.5cm, align=center] {\textbf{Fase di Propagazione}\\Replicazione e mutazione};
\node (att) [draw, rounded corners, fill=orange!10, right=of prop, text width=2.5cm, align=center] {\textbf{Fase di Attivazione}\\Il trigger scatta};
\node (exec) [draw, rounded corners, fill=red!10, right=of att, text width=2.5cm, align=center] {\textbf{Fase di Esecuzione}\\Payload espletato};

\draw [->, >=Stealth] (dorm) -- (prop);
\draw [->, >=Stealth] (prop) -- (att);
\draw [->, >=Stealth] (att) -- (exec);
\draw [->, >=Stealth, dashed, bend left=30] (dorm) to (att);
\end{tikzpicture}
\end{center}

Un virus tipicamente attraversa quattro fasi:
\begin{enumerate}
    \item \textbf{Fase dormiente:} il virus è inattivo. Verrà attivato da un evento (non tutti i virus hanno questa fase).
    \item \textbf{Fase di propagazione:} il virus inserisce una copia (talvolta mutata per sfuggire al rilevamento) di se stesso in altri programmi o aree di sistema. I cloni inizieranno a loro volta a propagarsi.
    \item \textbf{Fase di attivazione:} il virus viene attivato da un evento (es. conteggio di copie create) per svolgere la sua funzione.
    \item \textbf{Fase di esecuzione:} viene espletata la funzione (il payload).
\end{enumerate}

\subsection{Classificazione dei Virus}
\subsubsection{1. In Base al Target}
\begin{description}
    \item[Boot sector infector] Infetta il Master Boot Record (MBR) o il boot record, diffondendosi all'avvio.
    \item[File infector] Infetta i file considerati eseguibili dal sistema operativo.
    \item[Macro virus] Infetta file contenenti codice di scripting interpretato da un'applicazione (es. Office).
    \item[Virus multipartito] Infetta i file in modi diversi contemporaneamente. Per rimuoverlo, bisogna bonificare tutti i possibili luoghi di infezione.
\end{description}

\subsubsection{2. In Base alla Strategia di Camuffamento}
\begin{description}
    \item[Virus criptato] Utilizza la crittografia per occultarsi. Una porzione del virus genera una chiave casuale usata per cifrare il resto del codice, memorizzando la chiave col virus. Ad ogni replicazione, la chiave cambia, evitando firme costanti.
    \item[Virus furtivo (Stealth)] L'intero virus viene nascosto (es. tramite rootkit, compressione o mutazione) alterando le risposte del sistema per nascondere la propria presenza.
    \item[Virus polimorfo] Crea copie funzionalmente equivalenti ma con pattern di bit diversi per ogni copia. Spesso usa una \textit{mutation engine} (associata a crittografia).
    \item[Virus metamorfico] Si ridefinisce completamente a ogni iterazione tramite trasformazioni complesse. Cambia non solo l'aspetto, ma talvolta anche il comportamento, rendendo il rilevamento estremamente arduo.
\end{description}

\subsection{Macro Virus e di Scripting}

Un macro virus si attacca ai documenti e sfrutta le capacità di macro dell'applicazione (es. Office, PDF) per eseguirsi e propagarsi. 

\begin{warnbox}[Pericolosità dei Macro Virus]
\begin{itemize}
    \item \textbf{Indipendenza dalla piattaforma}: infettano contenuti trasversali; qualsiasi OS che supporti l'applicazione è vulnerabile.
    \item \textbf{Facilità di diffusione}: i documenti vengono condivisi molto più frequentemente rispetto agli eseguibili.
    \item \textbf{Controlli di accesso inefficaci}: poiché infettano documenti dell'utente, i controlli file system standard non possono prevenirli facilmente.
    \item Documenti come i PDF supportano componenti integrati (es. JavaScript), offrendo un'ampia superficie d'attacco.
\end{itemize}
\end{warnbox}

\section{I Worm}

Un worm è un programma che cerca attivamente altre macchine da infettare tramite vulnerabilità software o funzionalità di rete e, una volta insediatosi, le usa per sferrare attacchi ad altre macchine in rete.

\subsection{Meccanismi di Diffusione}
\begin{itemize}
    \item \textbf{E-mail / Messaggistica:} il worm si invia in allegato, eseguendosi all'apertura.
    \item \textbf{Condivisione file / USB:} si copia su supporti rimovibili, sfruttando l'esecuzione automatica o attendendo l'apertura manuale.
    \item \textbf{Esecuzione remota:} sfrutta vulnerabilità in servizi di rete.
    \item \textbf{Accesso remoto ai file / FTP:} si copia su altri sistemi sperando di essere eseguito.
    \item \textbf{Login remoto:} si autentica e copia/esegue se stesso via shell.
\end{itemize}
Come i virus, i worm presentano le fasi di dormienza, propagazione, attivazione ed esecuzione.

\subsection{Ricerca del Target (Scanning)}

Nella fase di propagazione, il worm scansiona la rete (\textit{fingerprinting}) per trovare nuovi host, creando una vasta rete distribuita. Le strategie di scanning includono:
\begin{description}
    \item[Casuale] Ogni host analizza IP casuali con un seed diverso.
    \item[Hit-list] L'attaccante precompila una lista di IP vulnerabili. Ogni macchina infetta riceve una porzione della lista da analizzare. Processo molto rapido.
    \item[Topologica] Utilizza le informazioni (tabelle di routing, cache ARP, rubriche) dell'host vittima per trovare nuovi bersagli.
    \item[Sottorete locale] Cerca target nella propria rete LAN, bypassando il firewall perimetrale.
\end{description}

\begin{notebox}[Andamento di propagazione di un Worm]
La propagazione avviene in tre fasi:
\begin{enumerate}
    \item Crescita esponenziale: molti host vulnerabili, infezione rapida.
    \item Crescita lineare: gli host sprecano tempo ad attaccare macchine già infette, il tasso rallenta.
    \item Coda lunga: la maggior parte è infetta; si cercano a fatica i restanti.
\end{enumerate}
\end{notebox}

\section{Vulnerabilità Lato Client}
\begin{description}
    \item[Drive-by-Download] Avviene quando un utente visita una pagina compromessa. Il codice sfrutta le vulnerabilità del browser per scaricare e installare malware \textbf{senza il consenso} dell'utente. A differenza dei worm, il malware attende passivamente la visita.
    \item[Watering Hole] Variante mirata del drive-by-download. Gli attaccanti infettano un sito web \textit{specifico} frequentemente visitato da un target aziendale o governativo.
    \item[Malvertising] Gli attaccanti acquistano spazi pubblicitari legittimi (Ad Network) che contengono script malevoli. Colpiscono siti affidabili senza doverli hackerare direttamente. Spesso mirati in base a fascia oraria o profilazione.
    \item[Clickjacking] L'interfaccia web viene manipolata. Attraverso \texttt{iframe} invisibili o trasparenti sovrapposti a bottoni apparentemente innocui, l'utente è ingannato a cliccare su comandi malevoli (es. autorizzare bonifici o permessi).
\end{description}

\section{Trojan}
\begin{defbox}[Trojan]
Un trojan è un programma in apparenza utile che contiene codice nascosto il quale, se invocato, svolge operazioni dannose, permettendo all'attaccante azioni indirette. Non si replicano autonomamente.
\end{defbox}

\newpage
I trojan possono comportarsi secondo tre modelli:
\begin{enumerate}
    \item \textbf{Esecuzione parallela}: Eseguono l'attività utile e il codice malevolo in parallelo.
    \item \textbf{Modifica funzioni}: Eseguono l'originale ma ne alterano parti (es. una shell di login che invia password, o comando \texttt{ps} che nasconde processi).
    \item \textbf{Sostituzione totale}: Un programma completamente fittizio.
\end{enumerate}

\section{Payload di un Malware}

Il payload è l'azione specifica che il malware esegue a beneficio dell'attaccante, spesso innescato da bombe logiche.

\subsection{Distruzione dei Dati}
\begin{warnbox}[Ransomware]
Cripta i dati dell'utente e richiede un riscatto (in criptovaluta) per la chiave di decrittazione.
\begin{itemize}
    \item \textbf{WannaCry (2017)}: Ha infettato migliaia di sistemi sfruttando una vulnerabilità Windows (EternalBlue), criptando file e chiedendo Bitcoin. Ha evidenziato l'importanza critica di backup offline e patch aggiornate.
\end{itemize}
\end{warnbox}

\subsection{Bombe Logiche}
Codice latente che "esplode" solo se si verificano precise condizioni (data, accesso utente specifico, assenza di ping a un server), cancellando dati o bloccando il sistema.

\section{Agenti di Attacco: Bot e Botnet}

\begin{defbox}[Bot e Botnet]
Un \textbf{bot} (o zombie) è un malware che assume segretamente il controllo di un computer. L'insieme coordinato di centinaia o migliaia di bot forma una \textbf{botnet}. 
\end{defbox}

\subsection{Uso delle Botnet}
Le botnet vengono impiegate per:
\begin{itemize}
    \item \textbf{Attacchi DDoS}: sovraccaricare bersagli con traffico massiccio.
    \item \textbf{Spamming}: invio massiccio di email, eludendo filtri anti-spam.
    \item \textbf{Sniffing e Keylogging}: rubare credenziali o dati sensibili (in chiaro tramite sniffer, o prima della cifratura tramite keylogger).
    \item \textbf{Diffusione malware}: usare i bot per distribuire worm o virus.
    \item \textbf{Adware e Click Fraud}: generare clic automatici su pubblicità.
\end{itemize}

\subsection{Controllo Remoto (Command and Control)}
La differenza principale tra un worm e un bot è la presenza di un'infrastruttura di Command and Control (CnC).
\begin{itemize}
    \item \textbf{Evoluzione}: In passato usavano canali IRC. Oggi usano HTTP/HTTPS o reti Peer-to-Peer per nascondere il traffico.
    \item \textbf{Protezione dei CnC}:
    \begin{itemize}
        \item \textit{DGA (Domain Generation Algorithms)}: Generano centinaia di domini fasulli al giorno per trovare il server CnC attivo, rendendo il blocco manuale inutile.
        \item \textit{Fast-flux DNS}: Si alternano rapidamente gli IP associati a un dominio tramite proxy nella botnet stessa.
    \end{itemize}
\end{itemize}

L'unica contromisura davvero efficace è l'identificazione e il \textbf{sequestro/disattivazione} dell'infrastruttura CnC (takedown).

\section{Furto di Credenziali}

\subsection{Keylogger e Spyware}
\begin{description}
    \item[Keylogger] Registra le digitazioni sulla tastiera. Poiché cattura i dati prima della cifratura, HTTPS è inutile in questo caso. I più evoluti filtrano i log per rubare solo le password.
    \item[Spyware e Trojan Bancari (Zeus)] Malware avanzati (come Zeus) che monitorano l'attività browser, reindirizzano traffico o iniettano form web fasulli (\textit{web injects}) per bypassare le contromisure grafiche (come le tastiere virtuali usate dalle banche).
\end{description}

\subsection{Phishing}
Ingegneria sociale via e-mail. L'attaccante invia mail mascherate da enti legittimi (es. banche) con senso d'urgenza, spingendo a cliccare su link fraudolenti per rubare le credenziali. Spesso inviato in massa via botnet.
\begin{exbox}[Spear-Phishing]
Variante altamente mirata (profilazione della vittima sui social) per colpire dirigenti o aziende specifiche, aumentando il tasso di successo (es. false autorizzazioni bancarie).
\end{exbox}

\section{Stealthing (Backdoor e Rootkit)}

Le tecniche di \textit{stealthing} nascondono la presenza del malware nel sistema.

\subsection{Backdoor}
Accesso segreto (trapdoor) per aggirare l'autenticazione. Possono essere legittime (\textit{maintenance hooks} per il debug del codice) o malevole (inserite da attaccanti).
Per prevenirle: revisione attenta del codice, penetration testing e monitoraggio di rete per anomalie (es. porte in ascolto impreviste).

\subsection{Rootkit}
\begin{defbox}[Rootkit]
Insieme di strumenti progettati per ottenere e mantenere l'accesso root al sistema in modo invisibile. Sovvertono le API del sistema operativo (es. comandi come \texttt{ls} o il Task Manager) per non mostrare file, processi o connessioni legate all'attaccante.
\end{defbox}

\subsubsection{Classificazione dei Rootkit}
\begin{description}
    \item[Persistente] Sopravvive al riavvio infettando chiavi di registro o file d'avvio. Più facili da rilevare offline.
    \item[Residente in memoria] Non lascia traccia sul disco; sparisce al riavvio ma elude gli antivirus.
    \item[User-mode] Gira nello spazio utente, intercetta e altera le chiamate alle librerie (es. \texttt{libc}).
    \item[Kernel-mode] Gira a livello kernel, modificando le strutture dati del kernel. Molto più potente e difficile da rilevare.
    \item[Virtual Machine Based] Installa un hypervisor malevolo (VMM) e fa girare l'intero OS vittima in una macchina virtuale invisibile all'OS stesso.
    \item[Modalità esterna / Firmware] Si nasconde nel BIOS, UEFI o firmware di periferiche, offrendo persistenza totale anche in caso di formattazione.
\end{description}
